\section{Conclusiones y trabajo futuro}
\label{sec:conclusiones-e4}

\subsection{Cierre del ciclo E1--E4}

El proyecto recorrió, a lo largo de cuatro entregas, el ciclo completo de construcción de un
sistema distribuido: análisis del dominio y prototipo de comunicación (E1), arquitectura de
microservicios sobre Docker Compose (E2), persistencia distribuida real con particionado y
tolerancia a fallos verificada experimentalmente, más un pipeline analítico paralelo (E3), y el
cierre integral de esta entrega —dos clientes completos, refactor a arquitectura en capas con
patrones de diseño reales, una pirámide de pruebas completa, un pipeline de CI/CD de siete
\emph{jobs}, observabilidad de las tres señales, y una evaluación experimental de calidad contra
ISO/IEC 25010. Ninguna pieza de las entregas anteriores se descartó: el mismo cluster
CockroachDB, el mismo esquema de datos particionado, y el mismo pipeline de Spark de la Entrega 3
siguen operando sin cambios detrás de la nueva capa de clientes y observabilidad. La
Tabla~\ref{tab:trazabilidad-e1-e4} hace explícita esta progresión, entrega por entrega, con la
evidencia concreta (ADR o sección del documento) que sustenta cada aporte.

\begin{table}[H]
\centering
\small
\caption{Trazabilidad del ciclo E1--E4: qué aportó cada entrega y dónde queda la evidencia.}
\label{tab:trazabilidad-e1-e4}
\begin{tabular}{@{}p{1.1cm}p{6.4cm}p{5.6cm}@{}}
\toprule
\textbf{Entrega} & \textbf{Aporte arquitectónico} & \textbf{Evidencia} \\
\midrule
E1 & Definición del dominio ISP y primer prototipo monolítico con comunicación por
     \emph{sockets}. & Sección~\ref{sec:introduccion} \\
E2 & Descomposición en cinco microservicios (\texttt{auth}, \texttt{ticket}, \texttt{ai},
     \texttt{notification}, \texttt{report}), REST + eventos Kafka, JWT/RBAC. &
     Sección~\ref{sec:introduccion} \\
E3 & Fragmentación horizontal sobre CockroachDB de 3 nodos con verificación formal, postura
     CAP/PACELC mixta, pipeline analítico paralelo con Spark. & ADR-0003, ADR-0004,
     Secciones~\ref{sec:diseno-esquema}, \ref{sec:pipeline-paralelo} \\
E4 & Refactor a arquitectura hexagonal con 6 patrones GoF, apps web y móvil nativas, canal de
     telemetría PE-U1 y correlación \texttt{CORREL} en Incidencias, pipeline de CI/CD de 7
     \emph{jobs}, observabilidad de 3 señales, evaluación experimental ISO/IEC 25010. &
     ADR-0005, ADR-0006, ADR-0007, ADR-0008, Secciones~\ref{sec:arquitectura-e4},
     \ref{sec:web}, \ref{sec:movil}, \ref{sec:pruebas-cicd}, \ref{sec:observabilidad},
     \ref{sec:iso25010} \\
\bottomrule
\end{tabular}
\end{table}

\emph{Nota sobre PE-U1}: pese a que su nombre remite a la Entrega 1, se verificó contra el
historial completo del repositorio que este canal de telemetría (sockets + gRPC + reloj de
Lamport) nunca se construyó en ninguna entrega anterior --- se construyó por primera vez en esta
entrega, junto con \texttt{CORREL} (ADR-0007, ADR-0008), y se documenta así explícitamente para
no sugerir una continuidad que no existió.

\subsection{Aporte de cada integrante}
\label{sec:aporte-integrantes}

La Tabla~\ref{tab:aporte-integrantes} declara, por área del repositorio, quién aportó la mayoría
de las líneas según el historial de control de versiones. El método es el guion
\texttt{scripts/aporte\_manuscrito.py}, que recorre \texttt{git log --numstat} sobre las rutas
indicadas, mapea los dieciséis nombres de autor del historial a los cuatro correos reales, y
descarta las dependencias y los artefactos generados: \texttt{node\_modules/}, \texttt{build/},
\texttt{dist/}, \texttt{target/}, los archivos de bloqueo de dependencias y los artefactos de
compilación de \LaTeX{} (\texttt{.bbl}, \texttt{.blg}, \texttt{.aux}, \texttt{.toc},
\texttt{.out}, \texttt{.log}). Esa exclusión importa: sobre \texttt{docs/latex/} son 581 de las
6.008 líneas insertadas, y sin descartarlas el aporte de C. Carpio al manuscrito pasa de un
8,38\,\% real a un 10,45\,\% aparente. El método declarado en este epígrafe y el método aplicado
en la tabla son el mismo guion, y cualquiera puede reproducir cada fila corriéndolo:
\texttt{python scripts/aporte\_manuscrito.py docs/latex/ --commit <commit>}.

\begin{table}[H]
\centering
\caption{Aporte por área, según líneas reales insertadas en el historial (dependencias y
artefactos generados excluidos)}
\label{tab:aporte-integrantes}
\small
\begin{tabular}{@{}p{4.2cm}p{9cm}@{}}
\toprule
\textbf{Área} & \textbf{Aporte principal} \\
\midrule
Backend (\texttt{services/}) & C. Pacheco 40,69\,\%, C. Carpio 27,45\,\%, R. Cando 25,01\,\%, J. Álvarez 6,85\,\% \\
Aplicación web (\texttt{apps/web/}) & C. Carpio 97,68\,\%, J. Álvarez 2,32\,\% \\
Aplicación móvil (\texttt{apps/mobile/}) & R. Cando 97,65\,\%, C. Carpio 2,35\,\% \\
Manuscrito (\texttt{docs/latex/}) & C. Pacheco 76,45\,\%, R. Cando 12,90\,\%, C. Carpio 8,38\,\%, J. Álvarez 2,27\,\% \\
Experimentos (\texttt{experimentos/}, \texttt{spark/}) & R. Cando 73,78\,\%, C. Pacheco 17,71\,\%, C. Carpio 5,20\,\%, J. Álvarez 3,31\,\% \\
Observabilidad (\texttt{ops/}) & C. Pacheco 100,00\,\% \\
CI/CD (\texttt{.github/}) & C. Carpio 73,39\,\%, J. Álvarez 13,39\,\%, C. Pacheco 6,94\,\%, R. Cando 6,29\,\% \\
Arquitectura (\texttt{docs/adr/}, \texttt{docs/diagrams/}) & \textbf{J. Álvarez 45,82\,\%}, C. Pacheco 22,52\,\%, R. Cando 16,28\,\%, C. Carpio 15,38\,\% \\
Contratos (\texttt{pacts/}, \texttt{docs/api/}) & \textbf{J. Álvarez 100,00\,\%} \\
\bottomrule
\end{tabular}

\vspace{2pt}
{\footnotesize\textit{Nota: cada fila se reproduce con
\texttt{python scripts/aporte\_manuscrito.py <rutas> --commit COMMIT\_ENTREGA}, tomada en el
commit de la entrega. Las líneas de dependencias y de artefactos generados quedan fuera del
denominador, como declara el epígrafe.}}
\end{table}

La distribución es desigual por área, no por persona: cada uno de los cuatro integrantes es autor
principal de al menos un área completa del repositorio. Jeremy Álvarez es, en particular, el autor
único del contrato de interfaz OpenAPI y de las pruebas de contrato asociadas, y el principal de
las tres vistas de arquitectura C4 —trabajo real y sustantivo que, sin embargo, no se tradujo en
líneas del manuscrito LaTeX: su aporte ahí es de una sola sección de la Entrega 3
(\texttt{conclusiones.tex}, julio) y esta misma tabla, escrita para esta entrega. Declararlo así,
en vez de omitirlo, es más útil que una cifra de commits que por sí sola no distingue entre
``no trabajó'' y ``trabajó en código y contratos, no en prosa''.

\subsection{Contribuciones concretas de esta entrega}

\begin{itemize}[nosep]
  \item Refactor verificable de \texttt{ticket-service} a cuatro capas con seis patrones GoF,
        cada uno documentado con el problema real que resolvía antes del cambio; el mismo patrón
        de capas se extendió a \texttt{auth-service} y \texttt{report-service} (Sección
        \ref{sec:hexagonal-e4}), con alcance declarado explícitamente para los tres servicios que
        quedaron fuera.
  \item Dos aplicaciones cliente completas y funcionales (web y móvil) consumiendo la misma API a
        través de un único \emph{gateway}.
  \item Contratos consumidor-proveedor verificados en verde contra el sistema real, no contra
        \emph{mocks} estáticos.
  \item Un pipeline de CI/CD de diez \emph{jobs} que actúa como compuerta de calidad real.
  \item Observabilidad de extremo a extremo con las tres señales correlacionadas por un
        identificador común.
  \item Una evaluación experimental honesta contra ISO/IEC 25010, que reporta tanto lo que
        cumple el umbral objetivo como lo que no, con causa raíz identificada en cada caso.
  \item Cobertura de pruebas de \texttt{ticket-service} en 81.1\,\% del código propio (por
        encima del objetivo del 70\,\%), priorizando con el propio reporte de JaCoCo las clases
        reales con mayor número de líneas sin cubrir.
\end{itemize}

\subsection{Trabajo futuro}

Directamente derivado de las brechas reales encontradas en esta misma entrega (Secciones
\ref{sec:pruebas-cicd}, \ref{sec:iso25010} y \ref{sec:discusion-e4}):

\begin{enumerate}[nosep]
  \item Obtener una licencia de CockroachDB (o migrar a CockroachDB Serverless) para levantar el
        límite de transacciones concurrentes que hoy degrada la eficiencia y la fiabilidad bajo
        carga real.
  \item Agregar protección contra fuerza bruta (\emph{rate-limiting} o bloqueo temporal de cuenta)
        en el inicio de sesión de \texttt{auth-service}.
  \item Ejecutar el protocolo de evaluación ISO 25010 a escala completa (10 repeticiones, 5 y 10
        minutos por escenario) una vez que el cronograma del equipo lo permita, siguiendo
        exactamente la metodología ya validada a escala reducida en esta entrega.
  \item Completar las pruebas instrumentadas de Android en un dispositivo/emulador real, y ampliar
        su cobertura más allá del flujo de inicio de sesión.
\end{enumerate}
